% =====================================================================
% Fiche de revision - Python Scripting - Lesson 4 : Django
% Compilation : pdfLaTeX (Overleaf), aucune option particuliere requise.
% =====================================================================

\documentclass[a4paper,10pt]{article}

\usepackage[utf8]{inputenc}
\usepackage[T1]{fontenc}
\usepackage[french]{babel}
\usepackage[a4paper,margin=1.3cm]{geometry}
\usepackage{xcolor}
\usepackage{multicol}
\usepackage{booktabs}
\usepackage{array}
\usepackage{tabularx}
\usepackage{enumitem}
\usepackage{listings}
\usepackage{titlesec}
\usepackage{tcolorbox}

% ---------------------------------------------------------------- couleurs
\definecolor{accent}{HTML}{1F5FA8}
\definecolor{codefg}{HTML}{A03030}
\definecolor{codebg}{HTML}{F4F6F8}
\definecolor{kw}{HTML}{0B5FA5}
\definecolor{cmt}{HTML}{6C8A54}
\definecolor{strcol}{HTML}{2E7D32}
\definecolor{danger}{HTML}{B3261E}

% ---------------------------------------------------------------- titres
\titleformat{\section}
  {\normalfont\large\bfseries\color{accent}}{}{0pt}{}
  [\vspace{-0.75em}{\color{accent}\rule{\linewidth}{0.7pt}}]
\titlespacing*{\section}{0pt}{1.4ex}{0.7ex}

\titleformat{\subsection}[runin]
  {\normalfont\bfseries}{}{0pt}{}[\;---\;]
\titlespacing*{\subsection}{0pt}{0.8ex}{0.4em}

% ---------------------------------------------------------------- code
\lstset{
  language=Python,
  basicstyle=\ttfamily\footnotesize,
  keywordstyle=\color{kw}\bfseries,
  commentstyle=\color{cmt}\itshape,
  stringstyle=\color{strcol},
  showstringspaces=false,
  columns=fullflexible,
  keepspaces=true,
  backgroundcolor=\color{codebg},
  frame=none,
  xleftmargin=0.3em,
  aboveskip=0.55em,
  belowskip=0.55em,
  breaklines=true,
}

% code en ligne
\newcommand{\co}[1]{\texttt{\textcolor{codefg}{#1}}}

% ---------------------------------------------------------------- boites
\newtcolorbox{piegebox}[1]{
  colback=red!3, colframe=danger, boxrule=0.7pt,
  left=5pt, right=5pt, top=4pt, bottom=4pt,
  title=#1, fonttitle=\bfseries\small, coltitle=white
}

\newtcolorbox{astucebox}[1]{
  colback=accent!5, colframe=accent, boxrule=0.7pt,
  left=5pt, right=5pt, top=4pt, bottom=4pt,
  title=#1, fonttitle=\bfseries\small, coltitle=white
}

\setlist[itemize]{leftmargin=1.1em, itemsep=0.15em, topsep=0.25em, parsep=0pt}
\setlength{\columnsep}{1.1em}
\setlength{\parindent}{0pt}
\renewcommand{\arraystretch}{1.15}


\begin{document}

% Desactive les espaces automatiques du francais avant : ; ! ?
% pour que le code en ligne (ex. "def f(request):") ne soit pas deforme.
\shorthandoff{;:!?}

\small

\begin{center}
  {\LARGE\bfseries\color{accent} Python Scripting --- Lesson 4}\\[0.3em]
  {\normalsize Django \quad$\cdot$\quad fiche de révision QCM}
\end{center}
\vspace{0.4em}

% =====================================================================
\begin{piegebox}{Les 16 réponses pièges à connaître par c\oe ur}
\begin{minipage}[t]{0.485\linewidth}
\begin{tabularx}{\linewidth}{@{}l>{\raggedright\arraybackslash}X@{}}
\toprule
\textbf{Situation} & \textbf{Résultat} \\
\midrule
\co{startproject} / \co{startapp} & \textbf{projet} / \textbf{application} \\
\co{makemigrations}       & écrit des fichiers, \textbf{ne touche pas} la base \\
\co{migrate}              & \textbf{seule} commande qui modifie la base \\
\co{sqlmigrate polls 0001}& \textbf{affiche} le SQL, n'exécute rien \\
app hors \co{INSTALLED\_APPS} & ses migrations sont \textbf{ignorées} \\
\co{Question(...)} sans \co{save()} & rien en base, \co{q.id} vaut \co{None} \\
2\textsuperscript{e} \co{save()}      & \co{UPDATE}, même \co{id} \\
\co{objects.all()} vide   & \co{<QuerySet []>}, pas \co{[]} \\
\bottomrule
\end{tabularx}
\end{minipage}\hfill
\begin{minipage}[t]{0.485\linewidth}
\begin{tabularx}{\linewidth}{@{}l>{\raggedright\arraybackslash}X@{}}
\toprule
\textbf{Situation} & \textbf{Résultat} \\
\midrule
\co{get()} sans résultat  & \co{DoesNotExist} (\co{filter()} : vide) \\
\co{level\_\_gt=5}        & \textbf{strictement} supérieur \\
\co{request.GET.get("p")} absent & \co{None} (\co{GET["p"]} : erreur) \\
\co{GET.get("min")}       & toujours une \textbf{chaîne} \co{'1'} \\
vue qui renvoie un \co{str} & erreur : il faut un \co{HttpResponse} \\
POST sans \co{csrf\_token} & \textbf{403} Forbidden \\
\co{\{\{ p \}\}} avec du HTML & \textbf{échappé} : affiché tel quel \\
\co{.all()} dans un template & \textbf{sans} parenthèses : \co{.all} \\
\bottomrule
\end{tabularx}
\end{minipage}
\end{piegebox}

\vspace{0.5em}
\begin{multicols}{2}

% =====================================================================
\section{Les commandes}

\begin{lstlisting}
python -m venv venv                     # Lesson 3
pip install Django==6.1.2
python -m django --version
django-admin startproject mysite        # projet
python manage.py runserver              # dev, :8000
python manage.py startapp polls         # application
python manage.py makemigrations polls   # ecrit
python manage.py sqlmigrate polls 0001  # montre le SQL
python manage.py migrate                # applique
python manage.py shell                  # console
python manage.py createsuperuser        # compte admin
\end{lstlisting}

\vspace{0.3em}
\co{runserver} est réservé au développement : en production, un serveur
WSGI/ASGI importe \co{wsgi.py} ou \co{asgi.py}.

% =====================================================================
\section{Projet et applications}

\begin{lstlisting}
mysite/
    manage.py        # script de gestion, on ne l'edite pas
    mysite/
        __init__.py  # fait du dossier un module
        settings.py  # configuration du projet
        urls.py      # routes du projet
        asgi.py      # point d'entree ASGI
        wsgi.py      # point d'entree WSGI
\end{lstlisting}

\begin{itemize}
  \item Un \textbf{projet} contient \textbf{plusieurs} applications.
  \item Une \textbf{application} est un paquet Python autonome et
        \textbf{réutilisable} (pages, admin, API, middlewares\dots).
  \item \co{admin}, \co{auth}, \co{sessions}\dots{} sont déjà des apps,
        visibles dans \co{INSTALLED\_APPS}.
  \item Activer la nôtre : \co{"polls.apps.PollsConfig"} (forme du cours)
        ou simplement \co{"polls"}.
\end{itemize}

% =====================================================================
\section{Vue et route}

\begin{lstlisting}
# polls/views.py
from django.http import HttpResponse

def index(request):
    return HttpResponse("Hello, world.")

# polls/urls.py
from django.urls import path
from . import views

urlpatterns = [
    path("", views.index, name="index"),
    path("hello", views.hello, name="hello"),
]

# mysite/urls.py
urlpatterns = [
    path("polls/", include("polls.urls")),
    path("admin/", admin.site.urls),
]
\end{lstlisting}

\subsection{URL finale} préfixe du projet + route de l'app :
\co{"polls/"} + \co{""} $\to$ \co{/polls/} ;
\co{"polls/"} + \co{"hello"} $\to$ \co{/polls/hello}.

\subsection{\co{include()}} délègue tout ce qui commence par le préfixe au
\co{urls.py} de l'application.

\subsection{\co{name=}} optionnel ; permet \co{\{\% url 'index' \%\}} et
\co{reverse("index")} sans écrire l'URL en dur.

\subsection{Ordre} Django teste les routes \textbf{dans l'ordre} et garde la
première qui correspond : une route \co{""} attrape-tout va \textbf{en
dernier} (cas de \co{theme\_material\_kit}).

\subsection{Slash final} \co{"hello"} ne répond pas à \co{/polls/hello/}
(404). \co{APPEND\_SLASH} ajoute un slash manquant, il n'en retire jamais.

% =====================================================================
\section{Migrations : les trois phases}

\begin{enumerate}[leftmargin=1.3em, itemsep=0.1em, topsep=0.2em]
  \item écrire ou modifier les modèles dans \co{models.py} ;
  \item \co{makemigrations polls} : fichier \co{0001\_initial.py} ;
  \item \co{migrate} : applique à la base (\co{db.sqlite3} par défaut).
\end{enumerate}

\begin{itemize}
  \item Comme Git : un historique de changements que l'on \emph{rejoue}
        pour reconstruire la base de production.
  \item Les migrations suivent la \textbf{structure}, pas les données.
  \item Premier \co{migrate} d'un projet vide : il crée déjà des tables,
        celles des apps préinstallées (\co{auth}, \co{admin}\dots).
  \item Changer de SGBD : dictionnaire \co{DATABASES} de \co{settings.py}.
\end{itemize}

% =====================================================================
\section{Modèles}

\begin{lstlisting}
from django.db import models

class Question(models.Model):
    question_text = models.CharField(max_length=200)
    pub_date = models.DateTimeField("date published")

    def __str__(self):
        return self.question_text

class Choice(models.Model):
    question = models.ForeignKey(Question,
                   on_delete=models.CASCADE)
    choice_text = models.CharField(max_length=200)
    votes = models.IntegerField(default=0)
\end{lstlisting}

\begin{itemize}
  \item Modèle = une \textbf{table} ; attribut = une \textbf{colonne}.
  \item Clé primaire \co{id} auto-incrémentée ajoutée \textbf{automatiquement}.
  \item \co{"date published"} : nom lisible (\emph{verbose name}), pas une
        valeur par défaut.
  \item \co{ForeignKey} : \emph{many-to-one}. Chaque \co{Choice} pointe vers
        \textbf{une} \co{Question}, une \co{Question} a \textbf{plusieurs}
        choix. En base : colonne \co{question\_id}.
  \item \co{CASCADE} : supprimer la question supprime ses choix.
        (\co{PROTECT} interdit, \co{SET\_NULL} met à vide.)
  \item \co{\_\_str\_\_} : change l'affichage (shell, admin), \textbf{pas}
        de migration. Sans lui : \co{Question object (1)}.
\end{itemize}

% =====================================================================
\section{Shell et ORM}

\begin{lstlisting}
>>> from polls.models import Choice, Question
>>> Question.objects.all()
<QuerySet []>
>>> from django.utils import timezone
>>> q = Question(question_text="What's new?",
...              pub_date=timezone.now())
>>> q.save()                # INSERT
>>> q.id
1
>>> q.question_text = "What's up?"
>>> q.save()                # UPDATE, id inchange
>>> Question.objects.all()
<QuerySet [<Question: What's up?>]>
>>> q.choice_set.create(choice_text="Choice 1", votes=2)
<Choice: Choice 1 (2 votes)>
\end{lstlisting}

\subsection{\co{choice\_set}} relation inverse créée par Django :
\co{<modèle en minuscules>\_set}. \co{create()} lie \textbf{et}
enregistre ; \co{count()} ne compte que les choix de \co{q}.

\subsection{Requêtes}
\begin{tabularx}{\linewidth}{@{}l>{\raggedright\arraybackslash}X@{}}
\toprule
\co{objects.all()}             & toutes les lignes (\co{QuerySet}) \\
\co{filter(level\_\_gt=5)}     & $>5$ ; \co{gte} $\geq$, \co{lt}, \co{lte} \\
\co{exclude(level\_\_lt=5)}    & le complément : $\geq 5$ \\
\co{get(id=1)}                 & \textbf{un} objet, sinon exception \\
\co{count()}                   & nombre de lignes \\
\bottomrule
\end{tabularx}

\vspace{0.3em}
\co{filter(level >= 5)} est une erreur : on écrit \co{champ\_\_lookup=valeur}.

\subsection{Fuseau horaire} avec \co{USE\_TZ = True}, les dates sont
stockées en \textbf{UTC} ; les templates les convertissent dans
\co{TIME\_ZONE}.

% =====================================================================
\section{Templates (MVC / MVT)}

Le traitement reste dans la vue, l'affichage passe dans un fichier HTML
placé dans \co{<app>/templates/}.

\begin{lstlisting}
from django.shortcuts import render

def showQuestionsWithTemplate(request):
    context = {"questions": Question.objects.all()}
    return render(request, "questions.html", context)
\end{lstlisting}

\begin{lstlisting}
<ul>
  {% for question in questions %}
    <li>{{ question.question_text }}
      <ul>
        {% for choice in question.choice_set.all %}
          <li>{{ choice.choice_text }} - {{ choice.votes }}</li>
        {% empty %}
          <li>No choice</li>
        {% endfor %}
      </ul>
    </li>
  {% endfor %}
</ul>
\end{lstlisting}

\begin{itemize}
  \item \co{\{\{ x \}\}} affiche ; \co{\{\% tag \%\}} exécute
        (\co{for}, \co{if}, \co{url}, \co{csrf\_token}) ;
        \co{\{\# \#\}} commente.
  \item Chaque clé du \co{context} devient une variable du template.
  \item Méthodes appelées \textbf{sans parenthèses} : \co{.all}, pas
        \co{.all()}.
  \item Blocs fermés explicitement (\co{endfor}, \co{endif}) :
        l'indentation ne compte pas.
  \item \textbf{Échappement automatique} : \co{<b>} s'affiche en texte. Un
        f-string dans \co{HttpResponse} ne protège rien
        $\to$ \co{django.utils.html.escape()}.
  \item \co{\{\% extends "base.html" \%\}} + \co{\{\% block content \%\}} :
        une mise en page commune pour toutes les pages.
\end{itemize}

% =====================================================================
\section{Paramètres de requête}

\co{?} ouvre la \emph{query string}, \co{\&} sépare les paramètres :
\begin{lstlisting}
/ynav/query/profile?first_name=John&last_name=Doe
\end{lstlisting}

\begin{lstlisting}
def query(request):
    p = request.GET.get("p")       # None si absent
    return render(request, "query.html", {"p": p})

low = int(request.GET.get("min"))  # str -> int
numbers = list(range(low, high + 1))   # max inclus
\end{lstlisting}

\co{request.GET["p"]} lève une erreur (sous-classe de \co{KeyError}) si le
paramètre manque : préférer \co{.get()}.

% =====================================================================
\section{Formulaires}

\begin{lstlisting}
from django import forms

class NameForm(forms.Form):
    name = forms.CharField(max_length=100)

def name_form(request):
    name = None
    if request.method == "POST":
        form = NameForm(request.POST)  # lie
        if form.is_valid():            # valide
            name = form.cleaned_data["name"]
    else:
        form = NameForm()              # vide
    return render(request, "form.html",
                  {"form": form, "name": name})
\end{lstlisting}

\begin{itemize}
  \item \co{forms.Form} : \textbf{aucune} table, aucune migration.
  \item \co{cleaned\_data} n'existe qu'\textbf{après} \co{is\_valid()}.
  \item \co{\{\% csrf\_token \%\}} obligatoire dans un \co{<form
        method="post">}, sinon \textbf{403}.
  \item \co{ModelForm} : champs générés depuis le modèle,
        \co{form.save()} insère la ligne, validateurs du modèle inclus.
  \item Après un POST réussi : \co{redirect(...)}
        (\emph{Post/Redirect/Get}), un F5 ne renvoie pas le formulaire.
\end{itemize}

% =====================================================================
\section{Admin}

\begin{itemize}
  \item Interface \textbf{CRUD} sur \co{/admin}, protégée par un login.
  \item Compte : \co{python manage.py createsuperuser} (le \emph{root}).
  \item Tables \emph{Users} et \emph{Groups} : fournies par l'app
        \co{auth}.
  \item Nos tables, dans \co{polls/admin.py} :
\end{itemize}
\begin{lstlisting}
admin.site.register([Question, Choice])
\end{lstlisting}

% =====================================================================
\section{Bibliothèques et comptes}

\subsection{Material Kit} \co{requirements.txt} + \co{pip install -r},
\co{"theme\_material\_kit"} dans \co{INSTALLED\_APPS} \textbf{avant} nos apps,
\co{migrate}, puis \co{path("", include("theme\_material\_kit.urls"))}
\textbf{en dernier}.

\subsection{Bulma} une feuille CSS ; \co{django-bulma-form-templates} pour
styler les formulaires.

\subsection{Auth} \co{django.contrib.auth.forms} fournit
\co{UserCreationForm} (inscription) et \co{AuthenticationForm} (connexion).
Un \emph{template tag} personnalisé (\co{<app>/templatetags/}) peut injecter
des classes CSS dans le rendu d'un champ.

% =====================================================================
\section{Avant de déployer}

\co{DEBUG = False}, \co{SECRET\_KEY} hors du code (variable d'environnement,
\co{django-environ}), \co{ALLOWED\_HOSTS} = les domaines du site, un vrai
serveur WSGI/ASGI à la place de \co{runserver}.

\end{multicols}

\vspace{0.3em}
\begin{astucebox}{Checklist des 60 dernières secondes}
\begin{minipage}[t]{0.485\linewidth}
\begin{itemize}
  \item \co{models.py} $\to$ \co{makemigrations} $\to$ \co{migrate}.
  \item Nouvelle app : \co{INSTALLED\_APPS} \textbf{et} \co{include()}.
  \item URL = préfixe projet + route app ; premier match gagnant.
  \item Une vue prend \co{request}, rend un \co{HttpResponse}.
\end{itemize}
\end{minipage}\hfill
\begin{minipage}[t]{0.485\linewidth}
\begin{itemize}
  \item \co{save()} : \co{INSERT} si pas d'\co{id}, sinon \co{UPDATE}.
  \item \co{gt} strict, \co{gte} large ; \co{get()} = un seul objet.
  \item Paramètres GET = chaînes ; \co{.get()} renvoie \co{None}.
  \item POST : \co{csrf\_token}, \co{is\_valid()}, puis \co{redirect}.
\end{itemize}
\end{minipage}
\end{astucebox}

\vspace{0.4em}
\begin{astucebox}{Vocabulaire des énoncés (le QCM est en anglais)}
\begin{minipage}[t]{0.485\linewidth}
\begin{itemize}
  \item \emph{view} : fonction qui reçoit la requête
  \item \emph{route / URLconf} : table des URLs
  \item \emph{query string / query parameter} : ce qui suit \co{?}
  \item \emph{migration} : changement de schéma enregistré
  \item \emph{QuerySet} : résultat paresseux d'une requête ORM
\end{itemize}
\end{minipage}\hfill
\begin{minipage}[t]{0.485\linewidth}
\begin{itemize}
  \item \emph{bound / unbound form} : avec ou sans données
  \item \emph{superuser} : administrateur total
  \item \emph{CRUD} : Create, Read, Update, Delete
  \item \emph{auto-escaping} : échappement automatique
  \item \emph{lookup} : suffixe \co{\_\_gt}, \co{\_\_gte}\dots
\end{itemize}
\end{minipage}
\end{astucebox}

\vspace{0.4em}
\begin{center}
\footnotesize\itshape
Format du contrôle : 6 propositions, 1 ou 2 bonnes réponses, une case fausse
annule la question.\\
Pour s'entraîner : \texttt{cd "Lesson 4" \&\& python quiz.py}
\end{center}

\end{document}
